\documentclass[11pt, oneside]{article}
\usepackage{geometry} 
\geometry{letterpaper}
\usepackage{graphicx}	

\usepackage{amssymb, amsmath}

\newcommand{\I}[2]{
  \ensuremath{ \mathbb{I}_{#1} \! \left[ #2 \right] }
}

\newcommand{\E}[2]{
  \ensuremath{ \mathbb{E}_{#1} \! \left[ #2 \right] }
}

\newcommand{\rnd}[2]{
  \ensuremath{ \frac{ \mathrm{d} #1}{ \mathrm{d} #2 } }
}

\begin{document}

\section{Decompositions}

$\mathbb{R}^{2}$, Lebesgue measures.

The joint probability density function
%
\begin{equation*}
p( x_{1}, x_{2} )
=
\frac{1}{ 2 \, \pi \, \sqrt{ \sigma_{1}^{2} \, \sigma_{2}^{2} \, (1 - \rho^{2} ) } }
\exp \left[
  -\frac{1}{2} Q( x_{1}, x_{2} )
\right]
\end{equation*}
%
where
%
\begin{equation*}
Q ( x_{1}, x_{2} )
=
\frac{ 
    \sigma_{2}^{2} \, x_{1}^{2} 
  - 2 \, \rho \, \sigma_{1} \, \sigma_{2} \, x_{1} \, x_{2}
  + \sigma_{1}^{2} \, x_{2}^{2}
}{
  \sigma_{1}^{2} \, \sigma_{2}^{2} \, (1 - \rho^{2} ).
}
\end{equation*}

The marginal density function over $X_{1}$ is given by integrating over $x_{2}$,
%
\begin{align*}
p ( x_{1} ) 
&=
\int \mathrm{d} x_{2} \, 
p( x_{1}, x_{2} )
\\
&=
\frac{1}{ 2 \, \pi \, \sqrt{ \sigma_{1}^{2} \, \sigma_{2}^{2} \, (1 - \rho^{2} ) } }
\int \mathrm{d} x_{2} \, 
\exp \left[
  -\frac{1}{2} Q( x_{1}, x_{2} )
\right].
\end{align*}

In order to evaluate this integral we need to isolate all of the $x_{2}$ terms 
in $Q ( x_{1}, x_{2} )$.  This requires completing the square for $x_{2}$,
%
\begin{align*}
\sigma_{2}^{2} \, x_{1}^{2} 
  - 2 \, \rho \, \sigma_{1} \, \sigma_{2} \, x_{1} \, x_{2}
  + \sigma_{1}^{2} \, x_{2}^{2}
&=
\sigma_{1}^{2} \,
\left(
  x_{2} - \rho \frac{ \sigma_{2} }{ \sigma_{1} } x_{1}
\right)^{2}
+ 
\sigma_{2}^{2} \, x_{1}^{2}
- \rho \, \sigma_{2}^{2} \, x_{1}^{2} 
\\
&=
\sigma_{1}^{2} \,
\left(
  x_{2} - \rho \frac{ \sigma_{2} }{ \sigma_{1} } x_{1}
\right)^{2}
+ 
\sigma_{2}^{2} \, x_{1}^{2} \, (1 - \rho^{2} ).
\end{align*}

Then
%
\begin{align*}
Q ( x_{1}, x_{2} )
&=
\frac{ 
    \sigma_{2}^{2} \, x_{1}^{2} 
  - 2 \, \rho \, \sigma_{1} \, \sigma_{2} \, x_{1} \, x_{2}
  + \sigma_{1}^{2} \, x_{2}^{2}
}{
  \sigma_{1}^{2} \, \sigma_{2}^{2} \, (1 - \rho^{2} )
}
\\
&=
\frac{ 
  \sigma_{1}^{2} \,
  \left(
    x_{2} - \rho \frac{ \sigma_{2} }{ \sigma_{1} } x_{1}
  \right)^{2}
  +   
  \sigma_{2}^{2} \, x_{1}^{2} \, (1 - \rho^{2} )
}{
  \sigma_{1}^{2} \, \sigma_{2}^{2} \, (1 - \rho^{2} )
}
\\
&=
\frac{ 
  \left(
    x_{2} - \rho \frac{ \sigma_{2} }{ \sigma_{1} } x_{1}
  \right)^{2}
}{
  \sigma_{2}^{2} \, (1 - \rho^{2} )
}
+
\frac{  x_{1}^{2} }{ \sigma_{1}^{2} }.
\end{align*}

Consequently, we can write the joint density function as
%
\begin{align*}
p( x_{1}, x_{2} )
&=
\frac{1}{ 2 \, \pi \, \sqrt{ \sigma_{1}^{2} \, \sigma_{2}^{2} \, (1 - \rho^{2} ) } }
\exp \left[
  -\frac{1}{2} Q( x_{1}, x_{2} )
\right]
\\
&=
\frac{1}{ 2 \, \pi \, \sqrt{ \sigma_{1}^{2} \, \sigma_{2}^{2} \, (1 - \rho^{2} ) } }
\exp \left[
  -\frac{1}{2} \left(
      \frac{ 
      \left(
        x_{2} - \rho \frac{ \sigma_{2} }{ \sigma_{1} } x_{1}
      \right)^{2}
    }{
      \sigma_{2}^{2} \, (1 - \rho^{2} )
    }
    +
    \frac{  x_{1}^{2} }{ \sigma_{1}^{2} }
  \right)
\right]
\\
&=
\frac{1}{ 2 \, \pi \, \sqrt{ \sigma_{1}^{2} \, \sigma_{2}^{2} \, (1 - \rho^{2} ) } }
\exp \Bigg[
  -\frac{1}{2}
      \frac{ 
      \left(
        x_{2} - \rho \frac{ \sigma_{2} }{ \sigma_{1} } x_{1}
      \right)^{2}
    }{
      \sigma_{2}^{2} \, (1 - \rho^{2} )
    }
\Bigg]
\exp \Bigg[
  -\frac{1}{2} \frac{  x_{1}^{2} }{ \sigma_{1}^{2} }
\Bigg].
\end{align*}

This factorization of the joint density function then allows us to reduce the 
marginal density function to a normal integral, which we can evaluate in 
closed form,
%
\begin{align*}
p ( x_{1} )
&=
\int \mathrm{d} x_{2} \, 
p( x_{1}, x_{2} )
\\
&=
\frac{1}{ 2 \, \pi \, \sqrt{ \sigma_{1}^{2} \, \sigma_{2}^{2} \, (1 - \rho^{2} ) } }
\int \mathrm{d} x_{2} \, 
\exp \Bigg[
  -\frac{1}{2}
    \frac{ 
      \left(
        x_{2} - \rho \frac{ \sigma_{2} }{ \sigma_{1} } x_{1}
      \right)^{2}
    }{
      \sigma_{2}^{2} \, (1 - \rho^{2} )
    }
\Bigg]
\exp \Bigg[
  -\frac{1}{2} \frac{  x_{1}^{2} }{ \sigma_{1}^{2} }
\Bigg]
\\
&=
\frac{1}{ 2 \, \pi \, \sqrt{ \sigma_{1}^{2} \, \sigma_{2}^{2} \, (1 - \rho^{2} ) } }
\exp \Bigg[
  -\frac{1}{2} \frac{  x_{1}^{2} }{ \sigma_{1}^{2} }
\Bigg]
\int \mathrm{d} x_{2} \, 
\exp \Bigg[
  -\frac{1}{2}
    \frac{ 
      \left(
        x_{2} - \rho \frac{ \sigma_{2} }{ \sigma_{1} } x_{1}
      \right)^{2}
    }{
      \sigma_{2}^{2} \, (1 - \rho^{2} )
    }
\Bigg]
\\
&=
\frac{1}{ 2 \, \pi \, \sqrt{ \sigma_{1}^{2} \, \sigma_{2}^{2} \, (1 - \rho^{2} ) } }
\exp \Bigg[
  -\frac{1}{2} \frac{  x_{1}^{2} }{ \sigma_{1}^{2} }
\Bigg]
\sqrt{ 2 \, \pi \, \sigma_{2}^{2} \, (1 - \rho^{2} ) }
\\
&=
\frac{1}{ \sqrt{2 \, \pi \, \sigma_{1}^{2} } }
\exp \Bigg[
  -\frac{1}{2} \frac{  x_{1}^{2} }{ \sigma_{1}^{2} }
\Bigg].
\end{align*}

After all of the that, the marginal density function is just a normal density function,
%
\begin{align*}
p ( x_{1} )
&=
\frac{1}{ \sqrt{2 \, \pi \, \sigma_{1}^{2} } }
\exp \Bigg[
  -\frac{1}{2} \frac{  x_{1}^{2} }{ \sigma_{1}^{2} }
\Bigg]
\\
&=
\text{normal} ( x_{1} \mid 0, \sigma_{1} ).
\end{align*}

We can use this same factorization of the joint density function to simplify the calculation of the
conditional density functions,
%
\begin{align*}
p ( x_{2} \mid x_{1} )
&=
\frac{ p ( x_{1},  x_{2} ) }{ p ( x_{1} ) }
\\
&=
\frac{
  \frac{1}{ 2 \, \pi \, \sqrt{ \sigma_{1}^{2} \, \sigma_{2}^{2} \, (1 - \rho^{2} ) } }
  \exp \Bigg[
    -\frac{1}{2}
        \frac{ 
        \left(
          x_{2} - \rho \frac{ \sigma_{2} }{ \sigma_{1} } x_{1}
        \right)^{2}
      }{
        \sigma_{2}^{2} \, (1 - \rho^{2} )
      }
  \Bigg]
  \exp \Bigg[
    -\frac{1}{2} \frac{  x_{1}^{2} }{ \sigma_{1}^{2} }
  \Bigg]
}{
  \frac{1}{ \sqrt{2 \, \pi \, \sigma_{1}^{2} } }
  \exp \Bigg[
    -\frac{1}{2} \frac{  x_{1}^{2} }{ \sigma_{1}^{2} }
  \Bigg]
}
\\
&=
\frac{1}{ \sqrt{ 2 \, \pi \, \sigma_{2}^{2} \, (1 - \rho^{2} ) } }
\exp \Bigg[
  -\frac{1}{2}
      \frac{ 
      \left(
        x_{2} - \rho \frac{ \sigma_{2} }{ \sigma_{1} } x_{1}
      \right)^{2}
    }{
      \sigma_{2}^{2} \, (1 - \rho^{2} )
    }
\Bigg].
\end{align*}

Miraculously, each conditional density function is also another normal density function,
%
\begin{align*}
p ( x_{2} \mid x_{1} )
&=
\frac{1}{ \sqrt{ 2 \, \pi \, \sigma_{2}^{2} \, (1 - \rho^{2} ) } }
\exp \Bigg[
  -\frac{1}{2}
      \frac{ 
      \left(
        x_{2} - \rho \frac{ \sigma_{2} }{ \sigma_{1} } x_{1}
      \right)^{2}
    }{
      \sigma_{2}^{2} \, (1 - \rho^{2} )
    }
\Bigg]
\\
&=
\text{normal} \left( 
  x_{2} \biggm| 
  \rho \frac{ \sigma_{2} }{ \sigma_{1} } x_{1}, 
  \sigma_{2} \sqrt{ 1 - \rho^{2} } 
\right).
\end{align*}

Because of the symmetric of the joint density function, the same calculations apply if 
when deriving the other marginal density function and corresponding conditional 
density functions.  The result just swaps all of the indices,
%
\begin{align*}
p ( x_{2} )
&=
\text{normal} ( x_{2} \mid 0, \sigma_{2} )
\\
p ( x_{1} \mid x_{2} )
&=
\text{normal} \left( 
  x_{1} \biggm| 
  \rho \frac{ \sigma_{1} }{ \sigma_{2} } x_{2}, 
  \sigma_{1} \sqrt{ 1 - \rho^{2} } 
\right).
\end{align*}

Finally, this is an example where some pattern matching can point to the marginal  
and conditional density functions directly.  Examining
%
\begin{equation*}
p( x_{1}, x_{2} )
=
\frac{1}{ 2 \, \pi \, \sqrt{ \sigma_{1}^{2} \, \sigma_{2}^{2} \, (1 - \rho^{2} ) } }
\exp \Bigg[
  -\frac{1}{2}
      \frac{ 
      \left(
        x_{2} - \rho \frac{ \sigma_{2} }{ \sigma_{1} } x_{1}
      \right)^{2}
    }{
      \sigma_{2}^{2} \, (1 - \rho^{2} )
    }
\Bigg]
\exp \Bigg[
  -\frac{1}{2} \frac{  x_{1}^{2} }{ \sigma_{1}^{2} }
\Bigg],
\end{equation*}
%
we might notice that both of the exponentiated quadratics define normal
density functions,
%
\begin{align*}
p( x_{1}, x_{2} )
&=
\frac{1}{ 2 \, \pi \, \sqrt{ \sigma_{1}^{2} \, \sigma_{2}^{2} \, (1 - \rho^{2} ) } }
\exp \Bigg[
  -\frac{1}{2}
      \frac{ 
      \left(
        x_{2} - \rho \frac{ \sigma_{2} }{ \sigma_{1} } x_{1}
      \right)^{2}
    }{
      \sigma_{2}^{2} \, (1 - \rho^{2} )
    }
\Bigg]
\exp \Bigg[
  -\frac{1}{2} \frac{  x_{1}^{2} }{ \sigma_{1}^{2} }
\Bigg]
\\
&=
\frac{1}{ \sqrt{ 2 \, \pi \, \sigma_{2}^{2} \, (1 - \rho^{2} ) } }
\exp \Bigg[
  -\frac{1}{2}
      \frac{ 
      \left(
        x_{2} - \rho \frac{ \sigma_{2} }{ \sigma_{1} } x_{1}
      \right)^{2}
    }{
      \sigma_{2}^{2} \, (1 - \rho^{2} )
    }
\Bigg]
\frac{1}{ \sqrt{ 2 \, \pi \, \sigma_{1}^{2} } }
\exp \Bigg[
  -\frac{1}{2} \frac{  x_{1}^{2} }{ \sigma_{1}^{2} }
\Bigg]
\\
&=
\text{normal} \left( 
  x_{1} \biggm| 
  \rho \frac{ \sigma_{1} }{ \sigma_{2} } x_{2}, 
  \sigma_{1} \sqrt{ 1 - \rho^{2} } 
\right) \,
\text{normal} ( x_{1} \mid 0, \sigma_{1} ).
\end{align*}

In other words, by isolating $x_{2}$ we have incidentally manipulated 
the joint density function into the desired conditional decomposition.
This approach isn't always useful, but it can save time when working
with similar calculations.

\section{Convolution 1}

Start with $X = \mathbb{R}$ and probability distribution defined by 
the normal density function
%
\begin{align*}
p( x )
&=
\text{normal} ( x \mid \mu, \sigma )
\\
&=
\frac{1}{ \sqrt{ 2 \, \pi \, \sigma^{2} } }
\exp \left[
  -\frac{1}{2} \frac{ ( x - \mu )^{2} }{ \sigma^{2} }
\right].
\end{align*}

We can left this probability distribution into a joint distribution 
over $X^{2}$ by introducing a conditional probability kernel defined
by the normal density functions
%
%
\begin{align*}
p( x' \mid x )
&=
\text{normal} ( x' \mid \eta - x, \tau )
\\
&=
\frac{1}{ \sqrt{ 2 \, \pi \, \tau^{2} } }
\exp \left[
  -\frac{1}{2} \frac{ ( x' - \eta + x )^{2} }{ \tau^{2} }
\right].
\end{align*}

The lifted joint density function is given by
%
\begin{align*}
p( x, x' )
&=
p( x' \mid x ) \, p( x )
\\
&=
\text{normal} ( x \mid \mu, \sigma ) \,
\text{normal} ( x' \mid \eta - x, \tau )
\\
&=
\frac{1}{ 2 \, \pi \, \sqrt{ \sigma^{2} \, \tau^{2} } }
\exp \left[
  -\frac{1}{2} Q(x, x')
\right],
\end{align*}
%
where
%
\begin{equation*}
Q(x, x') 
= 
  \frac{ ( x - \mu )^{2} }{ \sigma^{2} }
+ \frac{ ( x' - \eta + x )^{2} }{ \tau^{2} }.
\end{equation*}

The convolution of $p(x)$ with $p(x' \mid x)$ is given by 
projecting this joint density function to the auxiliary 
space,
%
\begin{align*}
p( x' ) 
&= 
\int \mathrm{d} x \, p( x, x' )
\\
&=
\frac{1}{ 2 \, \pi \, \sqrt{ \sigma^{2} \, \tau^{2} } }
\int \mathrm{d} x \, 
\exp \left[
  -\frac{1}{2} Q(x, x')
\right].
\end{align*}

To evaluate this integral, we'll need to isolate $x$ with some 
completion of some squares,
%
\begin{align*}
Q(x, x') 
&= 
  \frac{ ( x - \mu )^{2} }{ \sigma^{2} }
+ \frac{ ( x' - \eta + x )^{2} }{ \tau^{2} }
\\
&=
  \frac{ x^{2} - 2 \, \mu \, x + \mu^{2} }{ \sigma^{2} }
+ \frac{ x^{2} - 2 \, (x' - \eta) \, x + (x' - \eta)^{2} }{ \tau^{2} }
\\
&=
\frac{1}{ \sigma^{2} \, \tau^{2} } \bigg[
    ( \sigma^{2} + \tau^{2} ) \, x^{2}
  - 2 ( \tau^{2} \, \mu + \sigma^{2} \, (x' - \eta) ) \, x
  + \tau^{2} \, \mu^{2} + \sigma^{2} ( x' - \eta)^{2}
\bigg]
\\
&=
\frac{1}{ \sigma^{2} \, \tau^{2} } \bigg[ \quad
    ( \sigma^{2} + \tau^{2} ) \, \left( 
      x - \frac{  \tau^{2} \, \mu + \sigma^{2} \, (x' - \eta) }{ \sigma^{2} + \tau^{2} } 
    \right)^{2}
\\
&\hspace{16mm}
  + \tau^{2} \, \mu^{2} + \sigma^{2} ( x' - \eta)^{2}
  - \frac{ 
      \left( \tau^{2} \, \mu + \sigma^{2} \, (x' - \eta) \right)^{2}
    }{
      \sigma^{2} + \tau^{2} 
    }
\bigg]
\\
&=\quad
\frac{1}{ \sigma^{2} \, \tau^{2} } \bigg[
    ( \sigma^{2} + \tau^{2} ) \, \left( 
      x - \frac{  \tau^{2} \, \mu + \sigma^{2} \, (x' - \eta) }{ \sigma^{2} + \tau^{2} } 
    \right)^{2}
\bigg]
\\
&\quad+
\frac{1}{ \sigma^{2} \, \tau^{2} } \bigg[ \quad
    \frac{ 
       \tau^{4} \, \mu^{2} 
     + \sigma^{2} \, \tau^{2} \, \mu^{2} 
     + \sigma^{4} ( x' - \eta)^{2}
     + \sigma^{2} \, \tau^{2} \, ( x' - \eta)^{2}
    }{
      \sigma^{2} + \tau^{2} 
    }
\\
&\hspace{20mm}
  - \frac{ 
       \tau^{4} \, \mu^{2} 
     + 2 \tau^{2} \, \sigma^{2} \, \mu \, (x' - \eta)
     + \sigma^{4} (x' - \eta)^{2}
    }{
      \sigma^{2} + \tau^{2} 
    }
\hspace{15mm} \bigg]
\\
&=\quad
\frac{1}{ \sigma^{2} \, \tau^{2} } \bigg[
    ( \sigma^{2} + \tau^{2} ) \, \left( 
      x - \frac{  \tau^{2} \, \mu + \sigma^{2} \, (x' - \eta) }{ \sigma^{2} + \tau^{2} } 
    \right)^{2}
\bigg]
\\
&\quad+
\frac{1}{ \sigma^{2} \, \tau^{2} } \bigg[
    \frac{ 
       \sigma^{2} \, \tau^{2} \, \mu^{2} 
     - 2 \, \tau^{2} \, \sigma^{2} \, \mu \, (x' - \eta)
     + \sigma^{2} \, \tau^{2} \, ( x' - \eta)^{2}
    }{
      \sigma^{2} + \tau^{2} 
    }
\bigg]
\\
&=\quad
\frac{\sigma^{2} + \tau^{2} }{ \sigma^{2} \, \tau^{2} }
\left( 
  x - \frac{  \tau^{2} \, \mu + \sigma^{2} \, (x' - \eta) }{ \sigma^{2} + \tau^{2} } 
\right)^{2}
\\
&\quad+
\frac{ 
   \mu^{2} 
 - 2 \, \mu \, (x' - \eta)
 + ( x' - \eta)^{2}
}{
  \sigma^{2} + \tau^{2} 
}
\\
&=\quad
\frac{\sigma^{2} + \tau^{2} }{ \sigma^{2} \, \tau^{2} }
\left( 
  x - \frac{  \tau^{2} \, \mu + \sigma^{2} \, (x' - \eta) }{ \sigma^{2} + \tau^{2} } 
\right)^{2}
\\
&\quad+
\frac{ 
  ( x ' - (\mu + \eta) )^{2}
}{
  \sigma^{2} + \tau^{2} 
}.
\end{align*}

Now we can reduce the convolution operation to a normal integral,
%
\begin{align*}
p( x' ) 
&=
\frac{1}{ 2 \, \pi \, \sqrt{ \sigma^{2} \, \tau^{2} } }
\int \mathrm{d} x \, 
\exp \left[
  -\frac{1}{2} Q(x, x')
\right]
\\
&=
\frac{1}{ 2 \, \pi \, \sqrt{ \sigma^{2} \, \tau^{2} } }
\int \mathrm{d} x \, 
\exp \left[
  -\frac{1}{2} \frac{\sigma^{2} + \tau^{2} }{ \sigma^{2} \, \tau^{2} }
  \left( 
    x - \frac{  \tau^{2} \, \mu + \sigma^{2} \, (x' - \eta) }{ \sigma^{2} + \tau^{2} } 
  \right)^{2}
\right]
\exp \left[
  -\frac{1}{2} \frac{  ( x ' - (\mu + \eta) )^{2} }{ \sigma^{2} + \tau^{2} }
\right]
\\
&=
\frac{1}{ 2 \, \pi \, \sqrt{ \sigma^{2} \, \tau^{2} } } \,
\exp \left[
  -\frac{1}{2} \frac{  ( x ' - (\mu + \eta) )^{2} }{ \sigma^{2} + \tau^{2} }
\right]
\int \mathrm{d} x \, 
\exp \left[
  -\frac{1}{2} \frac{\sigma^{2} + \tau^{2} }{ \sigma^{2} \, \tau^{2} }
  \left( 
    x - \frac{  \tau^{2} \, \mu + \sigma^{2} \, (x' - \eta) }{ \sigma^{2} + \tau^{2} } 
  \right)^{2}
\right]
\\
&=
\frac{1}{ 2 \, \pi \, \sqrt{ \sigma^{2} \, \tau^{2} } } \,
\exp \left[
  -\frac{1}{2} \frac{  ( x ' - (\mu + \eta) )^{2} }{ \sigma^{2} + \tau^{2} }
\right]
\sqrt{2 \, \pi \, \frac{ \sigma^{2} \, \tau^{2} }{\sigma^{2} + \tau^{2} } }
\\
&=
\frac{1}{ \sqrt{ 2 \, \pi \, ( \sigma^{2} + \tau^{2} ) } } \,
\exp \left[
  -\frac{1}{2} \frac{  ( x ' - (\mu + \eta) )^{2} }{ \sigma^{2} + \tau^{2} }
\right]
\\
&=
\text{normal} \left( x' \mid \mu + \eta, \sqrt{ \sigma^{2} + \tau^{2} } \right).
\end{align*}

This has a neat interpretation.  The convolution operation takes in initial 
normal density function, translates it by $\eta$, and then inflates the standard 
deviation.

In the limit $\tau \rightarrow 0$, the convolution results in a pure 
translation,
%
\begin{equation*}
p( x' ) 
=
\text{normal} \left( x' \mid \mu + \eta, \sqrt{ \sigma^{2} + \tau^{2} } \right).
\end{equation*}
%
Specifically, this is exactly the pushforward of a normal density function
along the translation operation that we discussed in \textbf{Section Blah}.


At the same time, the convolution kernel in this limit acts like a singular 
Delta function that shifts an input $x$ to $x + \eta$.  This limit is 
also known as a linear convolution.

\section{Convolution 2}

Start with $X = \mathbb{I}$ and
%
\begin{equation*}
p( x ) 
= \text{Poisson}( x \mid \lambda )
= \frac{ \lambda^{x} e^{-\lambda} }{ x! }.
\end{equation*}

We can left this probability distribution into a joint distribution 
over $X^{2}$ by introducing a conditional probability kernel defined
by the conditional density functions
%
\begin{equation*}
p( x' \mid x )
=
\frac{ \delta^{x' - x} e^{-\delta} }{ (x' - x)! } I[ x \le x' ].
\end{equation*}

These conditional density functions look a bit awkward.  Let's verify 
that they're properly normalized for all $\delta$,
%
\begin{align*}
\sum_{x' = 0}^{ \infty }
p( x' \mid x )
&=
\sum_{x' = 0}^{ \infty }
\frac{ \delta^{x' - x} e^{-\delta} }{ (x' - x)! } I[ x \le x' ]
\\
&=
\sum_{x' = x}^{ \infty }
\frac{ \delta^{x' - x} e^{-\delta} }{ (x' - x)! }
\\
&=
\sum_{k = 0}^{ \infty }
\frac{ \delta^{k} e^{-\delta} }{ k! }
\\
&=
e^{-\delta} 
\sum_{k = 0}^{ \infty }
\frac{ \delta^{k} }{ k! }
\\
&=
e^{-\delta} 
e^{\delta}
\\
&=
1.
\end{align*}

With the validity of the convolution kernel verified, we can proceed
to the convolution itself,
%
\begin{align*}
p ( x' )
&=
\sum_{x = 0}^{ \infty }
p( x' \mid x ) \, p (x)
\\
&=
\sum_{x = 0}^{ \infty }
\frac{ \delta^{x' - x} e^{-\delta} }{ (x' - x)! } I[ x \le x' ] \,
\frac{ \lambda^{x} e^{-\lambda} }{ x! }
\\
&=
e^{-(\lambda + \delta)} 
\sum_{x = 0}^{ \infty }
I[ x \le x' ] \,
\frac{ 1 }{ x! \, (x' - x)! }
\delta^{x' - x} \, \lambda^{x}
\\
&=
\frac{ e^{-(\lambda + \delta)} }{ x'! }
\sum_{x = 0}^{ x' }
\frac{ x'! }{ x! \, (x' - x)! }
\delta^{x' - x} \, \lambda^{x}.
\end{align*}

This summation can be simplified by the binomial theorem,
%
\begin{equation*}
( a + b )^{K} = \sum_{k = 0}^{K} a^{k} \, b^{K - k}.
\end{equation*}
%
Consequently,
%
\begin{align*}
p ( x' )
&=
\frac{ e^{-(\lambda + \delta)} }{ x'! }
\sum_{x = 0}^{ x' }
\frac{ x'! }{ x! \, (x' - x)! }
\delta^{x' - x} \, \lambda^{x}
\\
&=
\frac{ e^{-(\lambda + \delta)} }{ x'! }
(\lambda + \delta)^{x'}.
\end{align*}

This, however, is just another Poisson density function,
%
\begin{equation*}
p ( x' )
=
\frac{ e^{-(\lambda + \delta)} }{ x'! }
(\lambda + \delta)^{x'}
=
\text{Poisson}( x' \mid \lambda + \delta ).
\end{equation*}
%
This particular convolution translates the intensity of the 
Poisson parameter.

\end{document}  